\section{Discussion and Future Work}
\label{app:extended-discussion-limitations}

This appendix expands the discussion and future-work map. It covers
selected-score scope, relation to activation SAEs, cross-family interpretation,
DLA composition, and possible extensions.

\subsection{What Sparse Readout Prism Decomposes}
\label{app:extended-discussion-score-scope}

\paragraph{Selected readout score.}
Each local score decomposition selects one scalar
\(s_q(h)=h^\top q\). This may be a selected vocabulary logit, a specified
pairwise logit difference, a group contrast, a vocabulary-mean contrast, or a
top-competitor margin. The hidden state's readout feature projection vector
\(\{h^\top d_i\}_i\) is available before selecting such a score; support for a
particular token or contrast becomes defined only after that vector is combined
with the direction coefficients \(\beta_i(q)\). Distribution-wide decompositions
therefore require measuring reconstruction over the relevant vocabulary slice or
token group. This distinction matters because a token can be supported above a
reference while still losing to a named alternative or a group contrast, and a
top-1 token can have only a small margin over top-ranked competitors. Reference
contrasts, pairwise logit differences, group contrasts, and top-competitor
margins therefore define distinct selected readout scores.

\paragraph{Local score decomposition.}
SRP provides local additive accounting for a selected readout score. For that
score, the displayed terms express the exact value as SAE feature terms plus
preprocessing and reconstruction error terms under the conventions in
Appendix~\ref{app:score-decomposition}. Contribution terms are reported with
those error terms, so the local score decomposition is tied to its reconstruction
quality.

\subsection{Relation to Activation SAEs}
\label{app:extended-discussion-activation-sae}

\paragraph{Fitted object.}
Sparse Readout Prism uses sparse autoencoder machinery, but its fitted object
differs from activation-SAE work. Activation SAEs learn sparse codes for a
distribution of hidden states
\citep{bricken2023monosemanticity,cunningham2023sparse,gao2024scaling}. Sparse
Readout Prism learns sparse codes for rows of the unembedding matrix \(\WU\).
The learned dictionary features are therefore readout directions used by token
rows, not hidden-state features discovered from model activations.

\paragraph{Context dependence.}
The hidden state enters after the readout SAE has been learned. A feature
contributes to a selected readout score when two quantities align: the selected
readout direction uses that SAE
feature through its sparse row coefficients, and the current hidden state has a
large dot product with the corresponding decoder direction. Hence the same
token row can have different dominant feature contributions across contexts:
the row coefficients are fixed, while the projection terms
\(h^\top d_i\) change with the hidden state.

\paragraph{Label status.}
SAE feature labels summarize tokens whose unembedding rows use a dictionary
feature strongly. They are row-level descriptors for the local score decomposition;
feature ids, signed contributions, and reconstruction errors remain the measured
objects of analysis.

\subsection{Model-Family Results}
\label{app:extended-discussion-cross-family}

\paragraph{Qwen evidence.}
Across Qwen sizes, the selected dictionary settings preserve held-out signed
readout scores and readout-distribution behavior on hidden states with reconstruction errors
small enough to support the reported local score decompositions. These rows provide
the fixed display setting and Qwen-family companion analyses used for most
qualitative figures.

\paragraph{Additional families.}
The additional model-family rows report the same metric suite under
different tokenizer, readout, and post-training conditions. They separate row
reconstruction, LM-head replacement, and selected-score reconstruction; the
paper reports all three before interpreting feature contributions.

\subsection{Feature-Resolved DLA}
\label{app:extended-discussion-prism-dla}

\paragraph{Readout versus components.}
Readout-side and component-side analyses answer different questions. The SRP
decomposition asks which unembedding-side SAE terms score a selected logit or contrast;
component-level attribution asks which residual-stream components project onto
the selected readout direction in the realized forward pass.

\paragraph{Residual terms.}
Combining the two views yields additive component-by-feature terms. The
analysis reports both the DLA additivity error for the selected readout score
and the readout SAE reconstruction error left outside the displayed feature
terms. Interventions are the natural next step for circuit-level claims.

\subsection{Future Work}
\label{app:extended-discussion-future-work}

\paragraph{Extensions.}
Natural extensions include controlled interventions on SAE decoder directions,
cross-model feature-alignment analyses, distributional decompositions of local
top-\(k\) frontiers, and broader model-family sweeps. The basis also
extends to the effective linear readouts used by tuned lenses,
logit-prism variants, and future-token lenses, so lens methods can be
compared by sparse signed feature terms as well as decoded tokens;
SRP-nominated terms give targets for causal patching, ablation, and
generation-level evaluation. A further extension factorizes both sides
of a selected margin at once, pairing an activation dictionary for
\(h\) with the readout SAE for \(q\), so that each term records
the alignment between an activation-side and a readout-side feature. All
depend on the reliability condition used in the main text: feature terms
are read when the selected readout score has a small reported residual.

\paragraph{Residual-stream census.}
SRP also measures what is readout-readable before a hidden state is reduced
to a single score. The projection vector \(p_i=h^\top d_i\) can be broad,
while the selected readout direction \(q\) selects the terms \(\beta_i(q)p_i\).
Future work could quantify projection mass used by the selected readout score, assigned to
competitors, or active but unused. The unused category is local to the
selected readout score; it may support alternative tokens, later positions, or other
computations. Comparing the split across layers, ambiguous prompts, and
incorrect generations could test whether models carry broad task state beyond
the immediate readout decision.
